\section{Stationary zeros of convex facet restrictions}
\label{app:three:facet-classification}

This appendix supplies the zero-set argument used in
\Cref{three:edge-classification}. A stationary facet zero may lie in the
facet interior or on one of its edges. The latter case is needed to
exclude a vanishing inward derivative at an edge contact.

\subsection{Subtracting an affine square}

\begin{lemma}\label{three:zero-span}
Let $q\in\mathcal P_3$. If every zero of $q$ lies in an affine plane
$\ell(x)=0$, with $\ell$ nonzero, then there is an $\epsilon>0$ such that
$q-\epsilon\ell^2\in\mathcal P_3$. Consequently an extreme ray of
$\mathcal P_3$ that is not an affine square has zeros affinely spanning
$\R^3$.
\end{lemma}
\begin{proof}
We use the rank-one subtraction criterion of
\citet[Lemma~4.3 in arXiv version~4]{Hildebrand2014MinimalZeros}:
for a copositive matrix $B$
and a nonzero vector $w$, there exists an $\epsilon>0$ with
$B-\epsilon ww^\top$ copositive if and only if $w^\top u=0$ for every
nonnegative zero $u$ of $B$.

Let $A\in\Sbb^4$ be the homogeneous coefficient matrix of $q$, and let
$T\in\R^{4\times8}$ have the columns $(1,v)$ for the cube vertices $v$.
The cone generated by $(1,x)$ for $x\in C_3$ equals $T\R_+^8$. Thus
$B=T^\top AT$ is copositive. Write $l\in\R^4$ for the homogeneous
coefficient vector of $\ell$ and set $w=T^\top l$. The matrix $T$ has
rank four, so $w\ne0$.

For a nonzero zero $u\ge0$ of $B$, put $s=\mathbf1^\top u>0$. There is
an $x\in C_3$ with $Tu=s(1,x)$, and
\[
 0=u^\top Bu=s^2q(x).
\]
Hence $w^\top u=l^\top Tu=s\ell(x)=0$. The criterion gives
$B-\epsilon ww^\top\in\COP_8$. Every $(1,x)$ has a nonnegative vertex
representation, so this says exactly that $q-\epsilon\ell^2\ge0$ on
$C_3$. For an extreme $q$, the decomposition
$q=(q-\epsilon\ell^2)+\epsilon\ell^2$ forces $q$ to be an affine square.
\end{proof}

\subsection{The stationary-facet lemma}

\begin{lemma}\label{three:stationary-facet}
Let $q\in\mathcal P_3$ have positive square coefficients, positive mixed
coefficients, and positive values at every cube vertex. Suppose a facet
restriction has a positive-semidefinite quadratic coefficient matrix and
a zero at which both facet derivatives vanish. If $q$ generates an
extreme ray, then it is an affine square.
\end{lemma}

\begin{proof}
Simultaneously complementing all three coordinates preserves the mixed
coefficients, so we may take the facet to be $z=0$. Its stationary zero
$u_0=(x_0,y_0)$ is not a corner, because all vertex values are positive.
It may be an edge point. We prove that all zeros of $q$ lie in an affine
plane, then apply \Cref{three:zero-span}.

\smallskip\noindent
\emph{Positive-definite facet block.}
Write
\begin{equation}\label{eq:three:pd-facet-form}
 q(u,z)=(u-u_0)^\top A(u-u_0)
       +2zv^\top(u-u_0)+\gamma z+\tau z^2,
\end{equation}
where
\[
 A=\begin{pmatrix}a&c\\c&b\end{pmatrix}\succ0,
 \qquad a,b,c>0,\qquad v=(r,s)^\top>0,\qquad \tau>0.
\]
The inward derivative at the zero gives $\gamma\ge0$. If the full
quadratic coefficient matrix is PSD, \eqref{eq:three:pd-facet-form} is a
sum of affine squares and $\gamma z$. Extremality and the positive square
coefficients make it a single affine square. Otherwise its Schur
complement
\[
 \kappa_0=\tau-v^\top A^{-1}v
\]
is negative.

For each $z\in[0,1]$, let $u(z)$ be the unique minimizer of $q(u,z)$ over
$[0,1]^2$, and put $F(z)=q(u(z),z)$. Thus $F\ge0$, $F(0)=0$, and every
zero of $q$ has the form $(u(z),z)$ with $F(z)=0$. We first describe this
minimizer path.

When both coordinates are free, its velocity is
\begin{equation}\label{eq:three:pd-velocity}
 x'(z)=\frac{cs-br}{ab-c^2},\qquad
 y'(z)=\frac{cr-as}{ab-c^2}.
\end{equation}
These velocities cannot both be positive. If both are nonpositive, one
coordinate reaches zero and the other then decreases to zero. For example,
on $x=0$ the free coordinate has velocity $y'=-s/b$, and the gradient in
the fixed coordinate increases with velocity $2(r-cs/b)\ge0$.
The lower bound therefore does not release.

Suppose instead that $x'>0$. Then $cs>br$ and $y'<0$. If $y$ reaches zero
first, $x$ changes to velocity $-r/a$, and the lower-bound gradient in $y$
increases with velocity $2(s-cr/a)>0$. The path ends at $00$. The strict
inequality follows from
\[
 cs>br,\quad ab>c^2
 \quad\Longrightarrow\quad s>\frac{br}{c}>\frac{cr}{a}.
\]
If $x$ reaches one first, $y$ continues with velocity $-s/b$. The
upper-bound gradient in $x$ decreases with velocity $2(r-cs/b)<0$, so
$x$ remains one until $y$ reaches zero. At $10$, the lower-bound gradient
in $y$ increases with velocity $2s$, while the upper-bound gradient in $x$
increases with velocity $2r$. The latter may reach zero and release $x$.
On the subsequent edge $y=0$, $x$ decreases with velocity $-r/a$ and $y$
stays zero because $s-cr/a>0$. Thus, up to relabeling, the possible paths
are
\begin{equation}\label{eq:three:pd-paths}
 \text{free}\ \longrightarrow\ x=0\ \longrightarrow\ 00,
 \qquad
 \text{free}\ \longrightarrow\ x=1\ \longrightarrow\ 10
       \ \longrightarrow\ y=0\ \longrightarrow\ 00.
\end{equation}
Some phases may be absent.

An initial $u_0$ on a facet edge gives a truncation of these paths. An
outward velocity at a lower bound immediately fixes that coordinate and
makes the other decrease; the lower-bound gradient then increases. An
outward velocity at an upper bound starts the upper-edge phase. An inward
velocity starts the free phase. Simultaneous transitions omit a phase.
If a free solution lies on a bound throughout a phase, either description
gives the same formulas. The symmetric case $y'>0$ follows by relabeling.
This accounts for every initial position allowed by the hypotheses.

The path is continuous and piecewise affine. The value function is
continuously differentiable and piecewise quadratic: on each phase the
free gradients vanish, and the fixed coordinates have zero velocity, so
$F'(z)=\partial_zq(u(z),z)$; the expressions agree at transitions.
The curvature on each phase is given in \Cref{tab:three:facet-curvature}.

\begin{table}[htbp]
\centering
\caption{Curvature of the two-coordinate minimum in
\eqref{eq:three:pd-facet-form}. The entries give $F''/2$.}
\label{tab:three:facet-curvature}
\begin{tabular}{ll}
\toprule
Active coordinates & Curvature\\
\midrule
Both free & $\kappa_0=\tau-v^\top A^{-1}v<0$\\
$x$ fixed, $y$ free & $\kappa_y=\tau-s^2/b$\\
$y$ fixed, $x$ free & $\kappa_x=\tau-r^2/a$\\
Both fixed & $\tau>0$\\
\bottomrule
\end{tabular}
\end{table}

For the upper-edge path in \eqref{eq:three:pd-paths},
$cs>br$ and $ab>c^2$ imply $r^2/a<s^2/b$, hence
$\kappa_x>\kappa_y$. Its curvature sequence is therefore
\begin{equation}\label{eq:three:pd-curvature-sequence}
 \kappa_0,\quad\kappa_y,\quad\tau,
       \quad\kappa_x,\quad\tau,
 \qquad \kappa_0<0,\quad\kappa_x>\kappa_y.
\end{equation}
The lower-edge path has the shorter sequence
$\kappa_0,\kappa_y,\tau$, or its relabeling.

These sequences imply that the zeros of $F$ for $z>0$ are at most two
isolated points or one interval. Here are the details. A negative-curvature
piece cannot have an interior zero of a nonnegative function. Nor can a
zero at an interior phase transition meet such a piece: differentiability
gives $F'=0$ at the zero, and the negative second derivative would give
negative values immediately on that side. A zero at $z=1$ can occur at
the end of a concave piece, but is then the final zero.

If $\kappa_y\ge0$, the tail starting with the first free-edge phase in
\eqref{eq:three:pd-curvature-sequence} is convex. If
$\kappa_y<0\le\kappa_x$, the tail starting with the first corner is
convex. In each case its nonnegative value function has a connected zero
set. If $\kappa_x<0$, both free-edge pieces are strictly concave. Only the
two corner pieces can have interior zeros, and each has at most one because
its curvature is positive. A terminal zero on a concave piece gives at
most one later zero after the first corner, again for a total of two.
The shorter path has the same simpler bound.

A zero interval requires one free-edge curvature to be zero. If
$\kappa_y=0$, then $\kappa_x>0$, and the following positive-curvature
corner leaves the interval with value and derivative zero and becomes
strictly positive. Every later curvature is positive. If $\kappa_x=0$,
the preceding corner has its unique zero at the entry to the interval;
the following corner is positive beyond its exit. Thus no separate zero
accompanies a zero interval. The argument for a lower-edge path is the
same. Across an interval, $u(z)$ is affine on one free-edge phase, so the
corresponding zeros of $q$ form a straight segment. The only zero on
$z=0$ is $(u_0,0)$. Its union with at most two points, or with one segment,
is contained in an affine plane. This proves the lemma for a PD facet
block.

\smallskip\noindent
\emph{Rank-one facet block.}
The positive diagonal and mixed coefficients give
$A=(d_1,d_2)^\top(d_1,d_2)$ with $d_1,d_2>0$, and the stationary zero
gives
\begin{equation}\label{eq:three:rankone-facet-form}
 q=(d_1x+d_2y-h)^2+2z(rx+sy+\eta)+\tau z^2,
 \qquad r,s,\tau>0.
\end{equation}
Its base zero set is the segment $d_1x+d_2y=h$, $z=0$.
Vertex positivity gives $0<h<d_1+d_2$ and $h\ne d_1,d_2$.
Put $\lambda_1=r/d_1$, $\lambda_2=s/d_2$, and first suppose
$\lambda_1<\lambda_2$.

For fixed $S=d_1x+d_2y$, minimizing $rx+sy$ uses coordinate $x$ first.
The minimum is
\[
 J(S)=
 \begin{cases}
  \lambda_1S,&0\le S\le d_1,\\
  \lambda_1d_1+\lambda_2(S-d_1),&d_1\le S\le d_1+d_2.
 \end{cases}
\]
For every $z>0$ this allocation is unique. The remaining minimum is
\begin{equation}\label{eq:three:allocation-value}
 F(z)=\min_{0\le S\le d_1+d_2}
       \bigl\{(S-h)^2+2z[J(S)+\eta]+\tau z^2\bigr\}.
\end{equation}
If $h<d_1$, its path is $y=0\longrightarrow00$. If $h>d_1$, it is
\[
 x=1\ \longrightarrow\ 10\ \longrightarrow\ y=0
       \ \longrightarrow\ 00,
\]
with weighted sums $h-\lambda_2z$, $d_1$, $h-\lambda_1z$, and $0$,
respectively. Its curvature sequence is
\begin{equation}\label{eq:three:allocation-curvature}
 \tau-\lambda_2^2,\quad\tau,
       \quad\tau-\lambda_1^2,\quad\tau,
 \qquad \tau-\lambda_1^2>\tau-\lambda_2^2.
\end{equation}
The preceding argument gives at most two isolated outside zeros or one
interval. Two isolated zeros could only be a first-corner zero followed
by a later zero after a negative-curvature phase where $x$ is free. If
that phase had nonnegative curvature, the tail would instead
be convex and its zero set connected.

This two-zero possibility is excluded by an inward derivative at the first
corner. At a zero $(1,0,z_A)$, vertex positivity gives $z_A\in(0,1)$.
The vertical edge restriction is a square, so, with $d_3=\sqrt\tau$,
\[
 z_A=\frac{h-d_1}{d_3}.
\]
The inward derivative in coordinate $x$ is nonnegative, giving
\[
 d_1(d_1-h)+rz_A\le0,
 \qquad\text{and hence}\qquad\lambda_1\le d_3.
\]
Thus $\tau-\lambda_1^2\ge0$, a contradiction to the required negative
curvature when $x$ is free. The base segment is therefore accompanied
by at most one outside point or one straight zero segment.

The base segment and a point lie in an affine plane. If an outside zero
segment lies on $y=0$, its restriction in \eqref{eq:three:rankone-facet-form}
must be an affine square; coefficient comparison gives
$r=d_1d_3$ and $\eta=-hd_3$. Its zeros satisfy $d_1x+d_3z=h$.
If it lies on $x=1$, the corresponding comparison gives
$s=d_2d_3$ and $r+\eta=-(h-d_1)d_3$, and its zeros satisfy
$d_2y+d_3z=h-d_1$. In either case the base segment and the outside segment
lie in the same plane
\[
 d_1x+d_2y+d_3z=h.
\]
If $\lambda_2<\lambda_1$, interchange the two coordinates.

It remains to consider $\lambda_1=\lambda_2=\lambda$. Now $q$ depends on
$x,y$ only through $S$, and its minimizing sum is
$S(z)=\max(h-\lambda z,0)$. On the free-sum phase,
\[
 F(z)=2(\eta+\lambda h)z+(\tau-\lambda^2)z^2.
\]
An interior positive-$z$ zero forces this polynomial to vanish identically:
nonnegativity makes its derivative zero there, and its factor $z$ leaves
a linear factor. Then $\eta=-\lambda h$, $\tau=\lambda^2$, and $q$ is
the affine square $(S+\lambda z-h)^2$. Otherwise a free-sum zero can occur
only at the final boundary $z=1$, and all corresponding points lie in the
plane $S+\lambda z=h$. After $S(z)$ reaches zero, the minimizer is $00$;
the resulting positive-curvature quadratic has at most one zero. A
positive-$z$ zero at an interior transition would likewise force the
preceding free-sum polynomial to vanish identically, because $F$ is
differentiable there. A transition at $z=1$ is a cube vertex and cannot
be a zero. A free-sum zero at $z=1$ leaves no later corner phase within the
domain. Thus, in the non-square case, the base segment is accompanied by
one point, or by one coplanar upper-facet segment. All zeros again lie in
an affine plane.

Both possible ranks of the PSD facet block have now been covered.
\Cref{three:zero-span} completes the proof.
\end{proof}
